\chapter{Démarche de développement et gestion de projet}
\label{chap:methodo}

\section{Une approche incrémentale}
Le projet a été conduit selon une \textbf{démarche incrémentale}, chaque incrément
produisant une version fonctionnelle et validée avant de passer au suivant. Cette approche a
permis de maîtriser la complexité croissante de la plateforme et de garantir, à tout moment,
un système opérationnel. Le tableau~\ref{tab:phases} synthétise les phases successives.

\begin{table}[h]
\centering\small
\caption{Phases de développement incrémentales.}
\label{tab:phases}
\begin{tabular}{>{\bfseries}p{2.0cm} p{4.2cm} p{7.0cm}}
\toprule
Phase & Thème & Livrables principaux \\
\midrule
Socle & Données et signal & Adaptateur unifié, pipeline de fenêtrage, extraction de features, préparation des cinq datasets \\
Modélisation & Apprentissage & Benchmark des cinq modèles, sélection par dataset, auto-encodeurs, CNN engrenages \\
Plateforme & Application & Backend FastAPI, frontend React, simulation temps réel (WebSocket), copilot \\
Industrialisation & Crédibilité & Persistance, authentification JWT, dérive, calibration, explicabilité, audit, versioning \\
P0 terrain & Déployabilité & Sévérité ISO 10816, ingestion universelle, baseline par machine, boucle GMAO, fausses alarmes, MQTT \\
P1 connectivité & Passage à l'échelle & Connecteur OPC-UA, copilot RAG, architecture time-series (TimescaleDB-ready) \\
\bottomrule
\end{tabular}
\end{table}

\section{Validation continue}
Chaque fonctionnalité a fait l'objet d'une \textbf{validation de bout en bout} avant
intégration : vérification des dimensions et distributions des données, contrôle des
métriques de benchmark, essais des routes de l'API, et tests fonctionnels des parcours
utilisateur. Plusieurs anomalies ont ainsi été détectées et corrigées au fil de l'eau
(chapitre~\ref{chap:defis}), notamment la fuite de données potentielle, le réentraînement
défaillant, la concurrence d'accès à la base et la stabilité des scores d'anomalie. Cette
discipline de validation explique la robustesse finale de la plateforme.

\section{Outils et environnement}
Le développement s'est appuyé sur un environnement Python isolé (environnement virtuel), une
gestion de version Git, et une organisation modulaire du code séparant clairement les
responsabilités : adaptateur de données, extraction de features, modèles, registre, routes
API, et interface. Les notebooks d'expérimentation (préparation des datasets, benchmark,
auto-encodeurs) sont conservés séparément du code de production, garantissant la traçabilité
des résultats et la reproductibilité des entraînements.

\section{Principes de qualité retenus}
Trois principes ont guidé l'ensemble du développement :
\begin{description}[leftmargin=1.5em, style=nextline]
  \item[Honnêteté des résultats.] Aucun chiffre n'est fabriqué : les performances sont
        mesurées sans fuite de données, et les indicateurs de fiabilité ne s'affichent qu'en
        présence de données réelles.
  \item[Robustesse avant fonctionnalité.] Une fonctionnalité n'est considérée comme acquise
        qu'une fois éprouvée face aux cas dégradés (données imparfaites, concurrence,
        erreurs).
  \item[Traçabilité.] Toute décision de l'IA est journalisée, tout modèle est versionné, et
        toute alerte peut être expliquée.
\end{description}
Ces principes, au-delà de leur intérêt académique, sont précisément ceux qu'exige un
déploiement industriel réel.
